\section{Extraction and quadratic normalization}
\label{sec:extraction}

The extraction and shortest-chain adjustment follow
\cite[\S\S12.3,12.5]{QED3Manuscript}. They solve different problems.
The extracted matrix becomes scalar only after summing over polymers.
The adjustment then converts the vanishing integral of each remaining
quadratic kernel into vanishing of the full kernel at zero connection.
This is the normalization used in the Cauchy estimate.

\subsection{Vacuum and mass extraction}
\label{subsec:rg-extraction}

For a balanced activity, its part of bidegree $(2,0)$ can be
written uniquely as
\[
 \mathscr E^{[2,0]}(X;\mathscr A,\mathscr C)
   =\int\dd x\,\dd x'\,
             \bar\psi(x)^T K_X(\mathscr A,\mathscr C;x,x')\psi(x').
\]
The matrix kernel is the ordinary antisymmetric kernel with
the first bar index one and the second bar index zero.
Antisymmetry accounts for the other ordering; in particular
\begin{equation}\label{eq:rg-bilinear-norm}
 \|\mathscr E^{[2,0]}(X;\mathscr A,\mathscr C)\|_{h_k}
   =h_{k,1}^2\sum_{\beta,\beta'}\int\dd x\,\dd x'
                     |K_X(\mathscr A,\mathscr C;x,x')_{\beta\beta'}|.
\end{equation}
The kernel is antiperiodic in each argument separately.

For a small $X$, choose one lift $\widetilde X$ and define
\begin{equation}\label{eq:rg-local-extractions}
 a_0(\mathscr E,X)
       =\operatorname{Vol}(X)^{-1}\mathscr E_{00}(X;0,0),\qquad
 a_2(\mathscr E,X)
       =\operatorname{Vol}(X)^{-1}
             \int_{\widetilde X}\dd x\int_{\widetilde X}\dd x'\,
                                                    K_X(0,0;x,x').
\end{equation}
Both cluster labels are averaged with measure $\eta_k^3/48$.
The second expression is independent of the lift: simultaneous
translation by a period multiplies its two arguments by the
same antiperiodic sign, whose product is one.

For an $M$-cube $\square$ put
\begin{equation}\label{eq:rg-counterterm-maps}
 \varepsilon_{\mathrm v}(\mathscr E,\square)
          =-\sum_{\substack{X\supset\square\\X\text{ small}}}
                                                     a_0(\mathscr E,X),
 \qquad
 M(\mathscr E,\square)
          =-\sum_{\substack{X\supset\square\\X\text{ small}}}
                                                     a_2(\mathscr E,X).
\end{equation}
Translations by $M\Z^3$ identify the small-polymer sets over
different cubes. They leave the two extracted quantities
unchanged. Thus the two sums do not depend on $\square$.
Write their values as $\varepsilon_{\mathrm v}(\mathscr E)$ and
$M(\mathscr E)$.

\begin{proposition}\label{prop:rg-scalar-counterterms}
There is a complex-linear scalar map $m(\mathscr E)$ such that
$M(\mathscr E)=m(\mathscr E)1$. Both $m$ and
$\varepsilon_{\mathrm v}$ are real on $K_k^x$, and
\begin{equation}\label{eq:rg-counterterm-bounds}
 |\varepsilon_{\mathrm v}(\mathscr E)|
                  \le K_0M^{-3}\|\mathscr E\|_k^x,
 \qquad
 |m(\mathscr E)|\le K_0M^{-3}e_k^{1/2}\|\mathscr E\|_k^x.
\end{equation}
They commute with restriction of the second connection argument
to zero.
\end{proposition}

\begin{proof}
Substitution of the generator transformations in a bilinear,
followed by uniqueness of its kernel, gives
\begin{equation}\label{eq:rg-extraction-matrix-symmetries}
 \begin{split}
 a_2(\mathscr E,\tilde gX)&=S a_2(\mathscr E,X)S^{-1},\\
 a_2(\mathscr E,X)&=\mathsf C^{-1}a_2(\mathscr E,X)^T\mathsf C,\\
 a_2(\mathscr E,X)&=\varkappa^{-1}
                         \overline{a_2(\mathscr E,X)}\varkappa
                                                   \quad(\mathscr E\in K_k^x).
 \end{split}
\end{equation}
For the first relation, the two variables and their labels
are transformed together, and the measure is unchanged.
For charge conjugation, moving the unbarred leg past the barred
leg cancels the minus sign in its transformation, and exchanges
the two kernel variables. For the last relation, the reality
map conjugates the coefficient and multiplies both legs by
$\varkappa$; its inverse transpose is $\varkappa^{-1}$.
The constant kernel has the corresponding scalar invariance
and is real for a real family.

Let $z$ be the center of $\square$ and take
$S=\gamma_1,\gamma_2,\gamma_3$ or $\gamma_5$.
The symmetry $(S,(1-r(S))z)$ fixes $\square$ and permutes the
small polymers containing it. The first relation in
\eqref{eq:rg-extraction-matrix-symmetries} therefore shows that
$M(\mathscr E)$ commutes with all four matrices. Their commutant
is scalar. In the fixed Dirac basis, $\gamma_3$ and
$\Sigma_3=-i\gamma_1\gamma_2\gamma_3$ are diagonal with the
four distinct pairs of diagonal entries
$(1,1),(-1,1),(-1,-1),(1,-1)$. A commuting matrix is diagonal;
commutation with $\gamma_1$ identifies its first two and last
two entries, and commutation with $\gamma_5$ identifies these
pairs. This proves the scalar assertion. The last relation in
\eqref{eq:rg-extraction-matrix-symmetries} makes this scalar
real for a real family.

The coefficient norm gives, for each small $X$,
\begin{equation}\label{eq:rg-single-extraction-bounds}
 \begin{split}
 \operatorname{Vol}(X)|a_0(\mathscr E,X)|
                    &\le\|\mathscr E(X;0,0)\|_{h_k},\\
 \operatorname{Vol}(X)\sum_{\beta,\beta'}
                              |a_2(\mathscr E,X)_{\beta\beta'}|
                    &\le h_{k,1}^{-2}
                               \|\mathscr E^{[2,0]}(X;0,0)\|_{h_k}.
 \end{split}
\end{equation}
Use $h_{k,1}^{-2}=e_k^{1/2}$ and
$\operatorname{Vol}(X)\ge M^3$. Each contribution is bounded
by the corresponding right side of
\eqref{eq:rg-counterterm-bounds}, times $\exp(-\kappa d_M(X))$.
The polymer sum over a fixed cube is at most $K_0$ by
$\kappa\ge\kappa_0$. For the mass use one diagonal entry of
the matrix sum, now known to be scalar. This proves the bounds.
Linearity follows from uniqueness of the antisymmetric kernels.
Restriction gives the same kernels at $(0,0)$, proving the
last assertion.
\end{proof}

Subtract these local terms only on small polymers:
\begin{equation}\label{eq:rg-normalized-remainder}
 (R\mathscr E)(X;\mathscr A,\mathscr C)
 =\mathscr E(X;\mathscr A,\mathscr C)
       -a_0(\mathscr E,X)\operatorname{Vol}(X)
       -\int_X\dd x\,\bar\psi(x)^Ta_2(\mathscr E,X)\psi(x).
\end{equation}
For large $X$ set $R\mathscr E=\mathscr E$. Then
\begin{equation}\label{eq:rg-exact-extraction-identity}
 \sum_X\mathscr E(X;\mathscr A,\mathscr C)
  =-\varepsilon_{\mathrm v}(\mathscr E)\ell^3
     -m(\mathscr E)\langle\bar\psi,\psi\rangle
                           +\sum_X(R\mathscr E)(X;\mathscr A,\mathscr C).
\end{equation}
Indeed, decompose the integral and volume of each small polymer
over its $M$-cubes, interchange the finite sums, and use
\eqref{eq:rg-counterterm-maps} on each cube. Every chamber
component receives the same scalar mass coefficient:
\[
 m(\mathscr E)\langle\bar\psi,\psi\rangle
 =m(\mathscr E)\frac{\eta_k^3}{48}
         \sum_c\sum_g\sum_{\beta=1}^4
                                      \bar\psi_\beta(c,g)\psi_\beta(c,g).
\]

The diagonal bilinear of a constant matrix $A$ on $X$ has
kernel $48\eta_k^{-3}1[x=x'\text{ over }X]A$ when both
arguments are taken over the same lift. Its integral is
$\operatorname{Vol}(X)A$, and its norm is
\[
 h_{k,1}^2\operatorname{Vol}(X)
                         \sum_{\beta,\beta'}|A_{\beta\beta'}|.
\]
These identities and \eqref{eq:rg-single-extraction-bounds}
show that
\begin{equation}\label{eq:rg-remainder-properties}
 \begin{split}
 (R\mathscr E)_{00}(X;0,0)&=0\quad(X\text{ small}),\\
 \int_{\widetilde X}\dd x\,\dd x'\,
                  (R\mathscr E)_{20}(X;0,0;x,x')&=0\quad(X\text{ small}),\\
 \|R\mathscr E\|_k^x&\le2\|\mathscr E\|_k^x.
 \end{split}
\end{equation}
For the bound, the two subtracted terms have distinct bidegrees
and together have norm at most $\|\mathscr E(X;0,0)\|_{h_k}$.
Apply the triangle inequality to
\eqref{eq:rg-normalized-remainder}, then the activity norm.

The map $R$ preserves the activity spaces and restriction to
zero. The subtracted kernels are independent of the connections,
and their two legs lie over $X$, so analyticity and locality
remain. A diagonal bilinear is invariant under every local
gauge. Its lattice, charge and reality transformations follow
from \eqref{eq:rg-extraction-matrix-symmetries}; the volume
and constant term have the same symmetries. This verifies
each requirement on the remainder, without assuming that
an individual matrix $a_2(\mathscr E,X)$ is scalar.

For later comparisons, the extraction has a precise local
meaning across different sides. Fix a lifted $M$-cube and
consider all lifted small polymers containing it. Their
projections are bijective onto the small polymers containing
its image, and their $d_M$ and volumes are unchanged.
If two families have the same constant and quadratic kernels
on every such lifted polymer at zero connections, then their
vacuum and mass matrices over this cube agree term by term in
\eqref{eq:rg-counterterm-maps}. Volume independence of a flow
will require this equality of kernels from the step comparison.
\subsection{Adjustment of the quadratic kernel}
\label{subsec:rg-adjustment}

The extraction makes the integral of the zero-field quadratic
kernel vanish. For contraction we need the kernel itself to
vanish. We achieve this by writing each ordinary leg as a
transported average and a sum of H\"older differences. This
changes the activity before evaluation and leaves its evaluated
value unchanged.

Let $X$ be small, choose a lift $\widetilde X$, and let
$x=(c,g)$ and $\tilde x=(\tilde c,\tilde g)$ lie over this lift.
Take all shortest face chains of its $M$-cubes between their
two cubes. For a chain $(\square_0,\ldots,\square_n)$ with
centers $z_i$, use the points
\[
 p_0=x,\qquad p_{i+1}=(z_i,g)\quad(0\le i\le n),\qquad
 p_{n+2}=\tilde x.
\]
The initial label is kept at each intermediate center.
Each distinct consecutive pair is admissible for $X$ of the
first kind. The paths in the first and last cube remain in
those cubes; between adjacent centers all six paths are the
straight segment across their common face. The last pair
can have different labels, and can lie inside one cluster.
For a coincident pair there is no difference term.

With $u_B(p,p')=U_B(p,p')$ and $u_B(p,p)=1$, set
\[
 U_i^\pm=\prod_{j<i}u_B(p_j,p_{j+1})^{\pm1}.
\]
The chain has no repeated cube, and
$n+1\le|X|_M\le c_X(L+1)$. Hence
\begin{equation}\label{eq:rg-chain-length}
 \sum_{i=0}^{n+1}|\pi p_i-\pi p_{i+1}|_1\le(n+3)M
          \le3c_XLM,\qquad
 \sum_{i:p_i\ne p_{i+1}}d_X(p_i,p_{i+1})^\alpha\le3c_XLM.
\end{equation}
The second estimate uses $d_X^\alpha\le\max(1,d_X)\le M$.
The transport bound is
$|u_B(p,p')^{\pm1}|\le\exp(e_k\|\operatorname{Im}B\|_\infty
|\pi p-\pi p'|_1)$, obtained by counting the bonds in each path.

For $B=\mathscr A+\mathscr C$, the domain decomposition and
$e_k^\varepsilon\le1/4$ give
$e_k\|\operatorname{Im}B\|_\infty\le e_k^\varepsilon/2$.
Assume in addition
\begin{equation}\label{eq:rg-adjustment-charge-condition}
 90c_XLM e_k^\varepsilon\le1.
\end{equation}
The modulus of every partial product $U_i^\pm$ is then at most
$\exp(1)$, uniformly in the chain and mesh.

Let $\operatorname{Mn}$ be the mean over the shortest chains.
Define the homomorphism $\Theta_{X,B}$ on ordinary generators by
\begin{equation}\label{eq:rg-chain-homomorphism}
 \begin{split}
 \Theta_{X,B}\psi(x)
  =\operatorname{Vol}(X)^{-1}\int_{\widetilde X}\dd\tilde x\,
    \operatorname{Mn}\left[
       U_{n+2}^+\psi(\tilde x)
       -\sum_{i:p_i\ne p_{i+1}}
           U_i^+d_X(p_i,p_{i+1})^\alpha\chi(p_i,p_{i+1})\right],\\
 \Theta_{X,B}\bar\psi(x)
  =\operatorname{Vol}(X)^{-1}\int_{\widetilde X}\dd\tilde x\,
    \operatorname{Mn}\left[
       U_{n+2}^-\bar\psi(\tilde x)
       -\sum_{i:p_i\ne p_{i+1}}
           U_i^-d_X(p_i,p_{i+1})^\alpha\bar\chi(p_i,p_{i+1})\right].
 \end{split}
\end{equation}
The spin components are unchanged. Both displayed images are
linear combinations of odd independent generators, so they
anticommute and determine a unique homomorphism on the
subalgebra of ordinary legs over $X$. Translation of the lift
by a period preserves every transport and chain and multiplies
each generator by the same antiperiodic sign. Thus the
homomorphism is independent of this choice of lift.

If the generators are evaluated at fields satisfying their
covariant H\"older relations, repeated substitution of
\[
 \psi(p_i)=u_B(p_i,p_{i+1})\psi(p_{i+1})
                    -d_X(p_i,p_{i+1})^\alpha\chi(p_i,p_{i+1})
\]
along a chain gives exactly the first expression in
\eqref{eq:rg-chain-homomorphism}; the barred relation uses
$u_B^{-1}$. Taking the two means therefore reproduces the
original leg. This proves equality after evaluation for
every polynomial in the ordinary legs over $X$.

The adjusted family is
\begin{equation}\label{eq:rg-adjustment-definition}
 \mathscr E^{\mathrm{nat}}(X)
   =\mathscr E(X)-\mathscr E^{[2,0]}(X)
           +\Theta_{X,\mathscr A+\mathscr C}
                                      \mathscr E^{[2,0]}(X)
                                      \qquad(X\text{ small}),
\end{equation}
and is unchanged on large polymers. The equality just proved
shows that the adjusted and original families have equal
evaluated values. Only bidegrees $(2,0),(1,1),(0,2)$ can change.

At zero connection all $U_i^\pm$ are one. The ordinary part
of an adjusted bilinear with kernel $K$ is therefore
\[
 \bar\psi_{\mathrm{av}}^T
    \left[\int_{\widetilde X}\dd x\,\dd x'\,K(x,x')\right]
                                                     \psi_{\mathrm{av}},
 \qquad
 \psi_{\mathrm{av}}=\operatorname{Vol}(X)^{-1}
                                        \int_{\widetilde X}\dd x\,\psi(x).
\]
Apply this to $R\mathscr E$, whose kernel integral vanishes.
With $E=(R\mathscr E)^{\mathrm{nat}}$, we obtain the stronger
normalization
\begin{equation}\label{eq:rg-adjusted-kernel-normalization}
 E_{00}(X;0,0)=0,\qquad E_{20}(X;0,0;\cdot,\cdot)=0
                                                     \quad(X\text{ small}).
\end{equation}

The same chain estimate controls the norm of adjustment.
Decompose $\Theta\psi=L_\psi+L_\chi$ by bidegree. For arbitrary
generator weights $(w_1,w_2)$,
\[
 \|L_\psi\|_w\le\exp(1)w_1,\qquad
 \|L_\chi\|_w\le\exp(1)(3c_XLM)w_2.
\]
The normalized site and label mean has total mass one.
Writing $s_K=\sum_{\beta,\beta'}\int|K_{\beta\beta'}|$ for
the balanced quadratic kernel and $F=\Theta\mathscr E^{[2,0]}$
gives the three bounds
\[
 \begin{aligned}
 \|F^{[2,0]}\|_w&\le\exp(2)w_1^2s_K,\\
 \|F^{[1,1]}\|_w&\le2\exp(2)(3c_XLM)w_1w_2s_K,\\
 \|F^{[0,2]}\|_w&\le\exp(2)(3c_XLM)^2w_2^2s_K.
 \end{aligned}
\]
Since $h_{k,2}=e_k^\varepsilon h_{k,1}$ and
$3c_XLM e_k^\varepsilon\le1/30$, summing these bounds yields
\begin{equation}\label{eq:rg-adjustment-norm}
 \|\mathscr E^{\mathrm{nat}}(X)\|_{h_k}
 \le\exp(2)(1+3c_XLM e_k^\varepsilon)^2\|\mathscr E(X)\|_{h_k}
                                               \le10\|\mathscr E(X)\|_{h_k}.
\end{equation}
Subtracting the original quadratic part before adding its
image uses additivity of the norm over bidegrees.

All new legs lie over $X$ or at its admissible chain pairs,
and all new transports read bonds with both ends in $X$.
Thus adjustment preserves locality. It preserves analyticity
because these transports are entire in the connection.
For the lattice covariance, symmetries carry shortest chains,
their centers and labels to the corresponding image chains.
For a real gauge the factors along a partial chain telescope
to the phase difference at its endpoints, giving the original
leg phase. Charge conjugation replaces $u_B$ by $u_B^{-1}$;
the reality map uses $\overline{u_{-\overline B}}=u_B$.
These identities on generators prove all four covariance
identities for $\Theta$, and hence the activity symmetries
for adjustment. Restriction to $\mathscr C=0$ commutes with
adjustment because it leaves the same total connection and
the same ordinary quadratic kernel. Therefore adjustment
preserves the complex and real activity spaces.